\scanpage{107}%
de Hodge), et par la condition que l'on a
\[
i_2=\overset{\handwritten{\rule{0.8em}{0.4pt}}}{i_1}\quad\text{i.e.}\quad i_2(\handwritten{\lambda})=\overset{\handwritten{\rule{2.2em}{0.4pt}}}{i_1(\handwritten{\bar{\lambda}})}\quad\text{pour tout}\quad\handwritten{\lambda\in}\underline{C}.
\]

Si à toute variété projective lisse $X$ sur $\underline{C}$ on associe sa cohomologie rationnelle $H^{\unsure{*|'}}(X,\underline{Q})=V$, de sorte que $V_{\underline{C}}=H^*(X,\underline{C})$ est isomorphe canoniquement (par la théorie de Hodge) à $H_{\mathrm{Hdg}}(X)=\handwritten{\coprod}H^q(X,\handwritten{\underline{\Omega}}^p_{X/\underline{C}})$, on voit qu'on trouve ainsi une pré-structure de Hodge sur $H^*(X,\underline{Q})$, d'où un \inserted{$\unsure{\otimes|\Omega}$-}foncteur naturel de la catégorie des motifs sur $\underline{C}$ dans la catégorie des pré-structures de Hodge. La conjecture de Hodge \unsure{stándart|standart} équivaut à dire que ce foncteur est \uline{pleinement fidèle}, i.e. que l\deleted{e}'homomorphisme naturel qui va du groupe de Galois de Hodge précédent $G$ dans le groupe de Galois motivique (associé au \inserted{$\otimes$-}foncteur de Betti $H_{\mathrm{Bet}}$) est un épimorphisme. \unsure{OU|Ou} encore que pour toute catégorie \inserted{$\underline{C}_o$} \inserted{de type fini} de motifs sur $\underline{C}$, désignant par $\underline{C}$ la \inserted{$\otimes$-}catégorie de pré-structures de Hodge engendrée par \deleted{\unsure{celle-ci,}} les $H_{\mathrm{Bet}}(M)$ pour $M\handwritten{\in}\operatorname{Ob}(\underline{C}_o)$, l'homomorphisme de groupes algébriques $G\handwritten{\longrightarrow}G_o$ associé au foncteur de Betti-Hodge $\underline{C}_o\handwritten{\longrightarrow}\underline{C}$ est un épimorphisme (i.e. surjectif sur les points à valeurs complexes, disons).
% A curved transposition mark takes “de type fini” from after “de motifs” to the caret before it.

On appelle \deleted{structure} \uline{polarisation} d'une pré-structure de Hodge de poids $n$ la donnée d'un accouplement de préstructures de Hodge
\[
\handwritten{\phi}:V\mathbin{\mathrm{x}}V\handwritten{\longrightarrow}\underline{Q}(n),
\]
où $\underline{Q}(n)$ est l'espace vectoriel trivial $\underline{Q}$ sur $\underline{Q}$, avec $\underline{Q}_{\underline{C}}=\underline{C}$ muni du \deleted{\unsure{la}} bidegré $(n,n)$, ayant la propriété que la forme hermitienne correspondante sur $V_{\underline{C}}$,
\margin{\handwritten{\uline{$\unsure{\psi|\phi}$}}}
\[
\handwritten{\psi}(x,y)=\handwritten{\phi}(x,\overset{\handwritten{\rule{0.5em}{0.4pt}}}{y})(\handwritten{i})^{p\handwritten{-q}}\quad\text{pour }x\text{ de bidegré }(p,\deleted{\unsure{n}}\handwritten{q}\unsure{-}p)
\]
soit définie positive. Une \uline{structure de Hodge} est une préstructure de Hodge admettant une polarisation. Les structures de Hodge forment une sous-$\otimes$-catégorie de la catégorie des pré-structures de Hodge. La théorie de Hodge nous assure que le foncteur de Betti-Hodge sur la catégorie des
% archive sheet number: 106
